\section{Introduction}
\label{sec:introduction}
A hot medium of strongly interacting, deconfined quarks and gluons, known as the quark-gluon plasma (QGP), is formed in ultra-relativistic heavy-ion collisions~\cite{Busza:2018rrf}. The multiplicity and pseudorapidity ($\eta$) distributions of charged particles produced in these collisions are critical observables for characterizing the initial conditions and the subsequent hydrodynamic evolution of the QGP~\cite{Romatschke:2017ejr}. The dependence of charged-particle multiplicity on the colliding system, center-of-mass energy, and collision geometry provides several insights, such as into nuclear shadowing and gluon saturation effects~\cite{Albacete:2014fwa}. It can further reveal the relative contributions to particle production from hard scattering and soft processes~\cite{Kharzeev:2000ph, dEnterria:2011twh}. Studying the charged hadron multiplicity and its dependence on $\eta$ is essential for understanding the formation and properties of the QGP in heavy-ion collisions. 

At the Relativistic Heavy Ion Collider (RHIC)~\cite{HARRISON2003235}, measurements of the system-size dependence of charged-particle $\eta$ density, denoted as $\dndeta$, have been performed for diverse collision systems, such as copper-copper, gold-gold (\auau), and deuteron-gold collisions at various center-of-mass energies~\cite{PHENIX:2000owy,PHENIX:2004vdg,PHENIX:2018hho,PHOBOS:2000wxz,PHOBOS:2003fjw,PHOBOS:2010eyu,BRAHMS:2001gci,BRAHMS:2001llo,STAR:2001eyo}. Similarly, the \alice, \atlas, and \cms experiments at the Large Hadron Collider (LHC) have reported $\dndeta$ at mid-rapidity for proton-proton, lead-lead (\pbpb), proton-lead, and xenon-xenon (\xexe) collisions at \tev energy scales~\cite{ALICE:2009wpl,ALICE:2010cin,ALICE:2010khr,ALICE:2010mty,ALICE:2015juo,ALICE:2015olq,ALICE:2015qqj,ALICE:2017pcy,ALICE:2018cpu,ALICE:2020swj,CMS:2010wcx,CMS:2010qvf,CMS:2010tjh, CMS:2011aqh,CMS:2014kix,CMS:2015zrm,CMS:2017shj,CMS:2019gzk,CMS:2024ykx,ATLAS:2011ag,ATLAS:2015hkr}. These measurements have revealed several key empirical trends:
\begin{enumerate*}[label=(\arabic*)]
\item charged-particle production approximately follows a power-law scaling with center-of-mass energy;
\item central heavy ion collisions show a steeper increase in $\langle\dndeta\rangle$ as a function of center-of-mass energy compared to proton-proton and proton-nucleus collisions;
\item the values of \dndeta, normalised by the number of participating nucleons, \npart, have a non-linear increase with \npart;
\item the shapes of the \npart dependence remain consistent across different collision energies.
\end{enumerate*}
These findings provide an opportunity to test scaling laws and models tuned to data from different energy regimes and evaluate their applicability to other collision systems.

This analysis uses data collected by \sphenix~\cite{PHENIX:2015siv, Belmont:2023fau} during the RHIC \auau run at a nucleon-nucleon center-of-mass energy of $\sqrtsnn=200\gev$, with a beam-beam crossing angle of 2 milliradians, taken in October 2024 under zero magnetic field conditions. The charged hadron yield per unit $\eta$, corrected for detector acceptance, reconstruction efficiency, combinatorial pairs, and contributions from secondary decays, is determined by counting tracklets within $|\eta|<1.1$. Tracklets are reconstructed by pairing clusters from the inner and outer layers of the Intermediate Silicon Tracker (\intt) that point back to the primary interaction vertex. Results are reported as a function of $\eta$ for different \auau centrality classes. To further explore particle production mechanisms, \dndeta values at mid-rapidity, scaled by number of participant nucleon pairs \npart/2, are studied as a function of \npart. These measurements are compared to previous RHIC results and to predictions from the heavy-ion event generators \hijing~\cite{Wang:1991hta}, \ampt~\cite{Lin:2004en}, and \epos~\cite{Pierog:2013ria}.